\chapter{Results and Validation}
\label{ch:results_and_validation}

This chapter presents the numerical outcomes of the validation pipeline when applied to $p_T$ spectra measured in high-energy collisions. The pipeline systematically evaluates single-component and multi-component models across ten multiplicity classes.

\section{Baseline Validations: BGBW and Tsallis Models}
To establish a rigorous baseline, we first fitted the spectra using the single-component BGBW and thermodynamically consistent Tsallis distributions.

The BGBW model yielded physically interpretable freeze-out parameters that perfectly align with expected literature trends (e.g., \citet{Khuntia2019, Rath2020}). At low multiplicity (class 21-30), the kinetic freeze-out temperature was found to be $T_{\text{kin}} = 132$ MeV with an average transverse flow velocity $\langle \beta \rangle = 0.31$. At high multiplicity (class 126-150), $T_{\text{kin}}$ decreased to $87$ MeV while the flow velocity increased to $\langle \beta \rangle = 0.66$. The fit quality was consistently excellent, with $\chi^2/ndf$ ranging strictly between $1.12$ and $2.06$ across all multiplicity bins.

The Tsallis baseline also performed well, though with slightly higher variance in fit quality ($\chi^2/ndf$ between $0.33$ and $14.29$).

\section{Evaluation of the Single-Component J\"uttner Model}
In contrast to the BGBW model, the single-component J\"uttner (Maxwell-Boltzmann) model completely failed to describe the data. The $\chi^2/ndf$ values ranged from $50.1$ at low multiplicity to an extreme $176.8$ at high multiplicity. This confirms that a simple exponential distribution without collective flow or hard scattering tails cannot adequately capture the $p_T$ spectra dynamics.

\section{Intrinsic Over-parameterization of the 2-Component J\"uttner Model}
The central claim of the evaluated manuscript proposed a multi-component model relying on moving J\"uttner distributions. Our validation pipeline tested a 2-component formulation of this model to evaluate its structural stability.

The \texttt{iminuit} optimizer failed to converge in 9 out of 10 multiplicity bins, returning invalid covariance matrices and Expected Distance to Minimum (EDM) values $> 0.1$. In a properly constrained model, the negative log-likelihood (or $\chi^2$) surface exhibits a well-defined parabolic minimum, allowing the Hessian matrix to be inverted to calculate parameter uncertainties. However, when a model is over-parameterized—as in the case of multiple unconstrained four-velocity vectors $U^\mu$—the likelihood manifold becomes extremely flat in certain directions. This manifests as highly correlated parameter pairs (with off-diagonal elements in the correlation matrix approaching $\pm 1.0$) and physically meaningless, diverging parameter uncertainties.

An exhaustive dense grid scan over the initial parameter space ($T$, $U_1$, $U_2$) confirmed that this convergence failure is not due to a poor initial guess (i.e., a local minimum trap), but rather an intrinsic topological feature of the over-parameterized model space. Even in isolated scenarios where the optimizer was forced to return a "success" state by artificially tightening parameter bounds, the resulting $\chi^2/ndf$ exceeded $172.5$. This indicates that even the "best" mathematical fit provides a physically unviable description of the experimental data.

We therefore definitively conclude that the proposed 2-component J\"uttner model is structurally unstable and unsuited for describing the multiplicity evolution of $p_T$ spectra.
